\documentclass[master.tex]{subfiles}

\begin{document}

% ======================== TITLE ================================
\ifSubfilesClassLoaded{%
  \begin{center}
    {\LARGE\bfseries\color{isublue}
     Chapter 4: Potential Outcomes and the Connection to Do-Calculus}\\[6pt]
    {\large Theory and Methods in Causal Inference}\\[4pt]
    {\normalsize Department of Statistics, Iowa State University}
  \end{center}
  \bigskip
  \hrule height 1.2pt \relax
  \smallskip
}{%
  \chapter{Potential Outcomes and the Connection to Do-Calculus}
}

% ======================== LEARNING OBJECTIVES ==================
\begin{tcolorbox}[defstyle, title={Learning Objectives}]
By the end of this chapter, students should be able to:
\begin{enumerate}
  \item Define the potential outcome $Y(t)$, state SUTVA precisely, and
        derive the consistency equation $Y = Y(T)$.
  \item Define the ATE and ATT, explain how they relate to each other,
        and re-express them as functionals of the structural equation for $Y$.
  \item Show that the linear SEM structural coefficient $\beta$ equals
        the ATE under homogeneous effects, and explain why ATE and ATT
        diverge under treatment effect heterogeneity.
  \item State the ignorability assumption, interpret it graphically as a
        back-door condition, and explain why overlap is a necessary
        companion assumption.
  \item Describe the SWIG construction of \citet{richardson2014ace},
        state the two SWIG-specific d-separation rules, and use them
        to derive the ignorability condition directly from the graph.
  \item Give a complete account of the Gaussian linear IV model in all
        three languages --- SEM, DAG/do-calculus, and potential outcomes
        --- and explain which language is most transparent for each task.
\end{enumerate}
\end{tcolorbox}

% ---------------------------------------------------------------
% Chapter-local styles and macros
% ---------------------------------------------------------------
% (\E defined in master.tex)
% (remarkstyle and remark environment are defined in causal_inference_master.tex)
% ==============================================================
\section{Motivation: A Third Language for Causality}
\label{w4:sec:motivation}
% ==============================================================

Part~I has built its identification theory in two languages: structural
equations (Chapter~1) and directed acyclic graphs with the do-calculus
(Chapters~2--3).  Both languages share a common syntax for the act of
intervention --- the operator $\doop(T{=}t)$ --- and both are
\emph{graphical} in the sense that causal assumptions are encoded as
directed edges.

A third language, the \emph{potential outcomes framework} of
\citet{neyman1923application} and \citet{rubin1974estimating}, approaches
causality differently.  Rather than modeling the mechanism by which $T$
affects $Y$, it defines the \emph{outcome that would be observed under each
possible treatment} as a random variable in its own right.  This language
dominates the biostatistics and econometrics literatures for one compelling
reason: it is the natural language for \emph{defining} causal estimands such
as the average treatment effect.

The strategic position of Chapter~4 in this course is deliberate.  Having
completed the full identification machinery of Chapters~1--3, students can now
see potential outcomes for what they are: a powerful vocabulary for
\emph{asking} causal questions that is, on its own, largely silent about
identification without additional structural assumptions.  The do-calculus
fills precisely that gap.

\begin{center}
\begin{tikzpicture}[
    stage/.style = {
      rectangle, rounded corners=5pt,
      draw=isublue, fill=defbg, thick,
      minimum width=2.8cm, minimum height=1.0cm,
      align=center, font=\small\bfseries\color{isublue}
    },
    hilight/.style = {
      rectangle, rounded corners=5pt,
      draw=accent, fill=defbg, thick,
      minimum width=2.8cm, minimum height=1.0cm,
      align=center, font=\small\bfseries\color{accent}
    },
    sub/.style  = {font=\scriptsize\itshape\color{darkgrey}, align=center},
    arr/.style  = {-{Stealth[length=7pt,width=5pt]},
                   line width=1.8pt, color=accent},
    node distance = 0.5cm
  ]
  \node[stage]   (ID)  {Identification\\(Chs.\ 1--3)};
  \node[hilight, right=1.2cm of ID]  (PO)
    {Potential Outcomes\\(Ch.\ 4) $\leftarrow$ \textit{here}};
  \node[stage, right=1.2cm of PO]   (RD)  {Research Designs\\(Chs.\ 5--8)};
  \node[stage, right=1.2cm of RD]   (EST) {Estimation\\(Part~III)};
  \draw[arr] (ID)  -- (PO);
  \draw[arr] (PO)  -- (RD);
  \draw[arr] (RD)  -- (EST);
  \node[sub, below=0.28cm of ID]  {DAGs, do-calculus,\\identification};
  \node[sub, below=0.28cm of PO]  {Define estimands;\\bridge frameworks};
  \node[sub, below=0.28cm of RD]  {Mediation, IV,\\propensity scores};
  \node[sub, below=0.28cm of EST] {Regression, IPW,\\doubly robust};
\end{tikzpicture}
\end{center}

\medskip
This chapter establishes the bridge between potential outcomes and the
do-calculus.  Section~\ref{w4:sec:po} defines the potential outcome framework
rigorously.  Section~\ref{w4:sec:estimands} develops the three principal causal
estimands: ATE and ATT.  Section~\ref{w4:sec:ignorability} connects
ignorability to the back-door criterion.  Section~\ref{w4:sec:swig} introduces
Single World Intervention Graphs (SWIGs) as a unified graphical framework.
Section~\ref{w4:sec:synthesis} revisits the Gaussian IV model in all three
languages to complete the Part~I synthesis.

\subsection{Three Languages for the Same Intervention: The Causal Trinity}
\label{w4:subsec:trinity}

The following table was first introduced in Chapter~1 and is reproduced
here for reference.  Throughout this course we will use these three
complementary languages --- collectively the \emph{causal trinity} ---
for expressing causal questions.

\begin{center}
\renewcommand{\arraystretch}{1.45}
\begin{tabular}{@{\hspace{6pt}}p{4.2cm}@{\hspace{10pt}}p{7.8cm}@{\hspace{6pt}}}
\toprule
\textbf{Framework} & \textbf{How the intervention $T{=}t$ is represented} \\
\midrule
Structural equations (SEM)
  & Replace the equation for $T$ with $T := t$ \\
DAG / do-calculus
  & Graph surgery producing $f(y \mid \doop(T{=}t))$ \\
Potential outcomes
  & Define the counterfactual variable $Y(t)$ \\
\bottomrule
\end{tabular}
\end{center}

All three members of the causal trinity describe the same conceptual
operation: externally setting the treatment to $t$.  Structural equations
describe the data-generating mechanism; DAGs make conditional independence
and identification transparent; and potential outcomes provide the most
natural notation for defining causal estimands such as the ATE and ATT.

\subsection{Part I Roadmap: Where Chapter 4 Sits}
\label{w4:subsec:roadmap}

Chapters~1--3 have built the identification machinery: SEMs define what
intervention means (Chapter~1), DAGs encode assumptions and license
d-separation (Chapter~2), and the do-calculus turns graph surgery into
identification formulas (Chapter~3).  What the first three chapters
\emph{do not} provide is a language for naming the quantities being
identified.  That is the job of potential outcomes.

\begin{center}
\renewcommand{\arraystretch}{1.45}
\begin{tabular}{@{\hspace{4pt}}c@{\hspace{10pt}}p{3.8cm}@{\hspace{10pt}}p{6.8cm}@{\hspace{4pt}}}
\toprule
\textbf{Chapter} & \textbf{Framework} & \textbf{New contribution} \\
\midrule
1--3 & SEM, DAGs, do-calculus
  & Identification: express $f(y \mid \doop(t))$ as a
    functional of observable distributions \\
4 & Potential outcomes
  & Definition: name the causal estimands
    $\tau_{\mathrm{ATE}} = \E[Y(1)-Y(0)]$,
    $\tau_{\mathrm{ATT}}$; connect them to $f(y \mid \doop(t))$
    and to the structural equation $g(T, \mathbf{X}, U_Y)$ \\
\bottomrule
\end{tabular}
\end{center}

The division of labor is therefore:
\[
  \textbf{Define the causal quantity}
  \quad\text{(potential outcomes)}
  \qquad
  \textbf{Identify it from data}
  \quad\text{(DAGs and do-calculus).}
\]

% ==============================================================
\section{The Neyman--Rubin Potential Outcomes Framework}
\label{w4:sec:po}
% ==============================================================

\subsection{The Potential Outcome}

\begin{defn}{Potential Outcome \citep{neyman1923application,rubin1974estimating}}
\label{w4:defn:po}
For each unit $i$ and each possible treatment value $t \in \mathcal{T}$,
the \emph{potential outcome} $Y_i(t)$ is the value of the outcome that
unit $i$ \emph{would have exhibited} had its treatment been set to $t$,
possibly contrary to fact.

The collection $\{Y_i(t) : t \in \mathcal{T}\}$ is called the
\emph{schedule of potential outcomes} for unit $i$.  In the binary case
$\mathcal{T} = \{0,1\}$, the schedule is $(Y_i(0),\, Y_i(1))$.
\end{defn}

\paragraph{Interpretation.}
$Y_i(1)$ is the outcome unit $i$ would achieve \emph{under treatment};
$Y_i(0)$ is the outcome under \emph{control}.  Only one of these is ever
observed for any given unit.  The unobserved potential outcome is called the
\emph{counterfactual}.

\begin{warn}{The Fundamental Problem of Causal Inference}
\citep{holland1986statistics}\ For any unit $i$, at most one potential
outcome $Y_i(t)$ is observed.
The individual-level causal effect $Y_i(1) - Y_i(0)$ is therefore
\emph{never directly observable}.  Causal inference is an inference problem
precisely because this quantity must be recovered from a population of units
rather than from a single unit.
\end{warn}

\subsection{SUTVA and Consistency}

The potential outcome $Y_i(t)$ is only well-defined if it does not depend
on the treatments assigned to \emph{other} units, and if there is a single,
unambiguous version of each treatment level.  These requirements are
formalized as SUTVA.

\begin{defn}{SUTVA \citep{rubin1980randomization}}
\label{w4:defn:sutva}
The \emph{stable unit treatment value assumption} has two components:
\begin{enumerate}
  \item \textbf{No interference.}  The potential outcome $Y_i(t)$ depends
        only on unit $i$'s own treatment, not on the treatments of other
        units: $Y_i(t_1, \ldots, t_n) = Y_i(t_i)$.
  \item \textbf{No hidden versions.}  For each treatment level $t$, there
        is a single well-defined version.  If two units both receive $T=1$,
        the treatment is the same in both cases.
\end{enumerate}
\end{defn}

\noindent\textbf{The consistency equation.}
Under SUTVA, the observed outcome $Y_i$ equals the potential outcome at the
treatment actually received:
\begin{equation}
  Y_i \;=\; Y_i(T_i).
  \label{w4:eq:consistency}
\end{equation}
This equation is the bridge between the potential outcome world and the
observed data world.  It fails whenever SUTVA fails: if treatment spills
over between units (e.g., vaccination provides herd immunity), or if the
treatment label $T=1$ covers multiple distinct interventions.

\subsection{Connection to the Do-Operator}

Under SUTVA and the structural assumptions of Chapters~1--3, the potential
outcome and the interventional distribution are linked by:

\begin{prop}{Potential Outcomes and the Do-Operator}
\label{w4:prop:po-do}
Under SUTVA,
\begin{equation}
  Y(t) \;\sim\; f\!\left(y \mid \doop(T{=}t)\right).
  \label{w4:eq:po-do}
\end{equation}
Equivalently, $\E[Y(t)] = \E[Y \mid \doop(T{=}t)]$.
\end{prop}

\noindent\textit{Proof sketch.}
In the mutilated SEM $\Gcal_{\overline{T}}$, the structural equation for
$T$ is replaced by $T := t$ for every unit.  The resulting value of $Y$ for
unit $i$ is determined solely by $t$ and unit $i$'s own background variables
--- which is precisely what the potential outcomes framework defines as
$Y_i(t)$.  SUTVA's no-interference condition ensures $Y_i(t)$ does not
depend on other units' treatment values, so the marginal distribution of $Y$
in $\Gcal_{\overline{T}}$ equals the marginal distribution of $Y(t)$.  By
definition of the do-operator, the former is $f(y \mid \doop(T{=}t))$.
\hfill$\square$

\begin{remark}
Equation~\eqref{w4:eq:po-do} shows that the do-operator and the potential
outcome notation are two ways of writing the same quantity.  The
do-operator makes the \emph{interventional/observational distinction}
syntactically explicit: $f(y \mid \doop(t)) \neq f(y \mid T{=}t)$
whenever $T$ is endogenous (established formally in Chapter~1).
The potential outcome notation $Y(t)$ is silent on this distinction ---
it does not carry a symbol that alerts the reader to the gap.  This is the
do-calculus language's primary advantage for identification.
\end{remark}

% ==============================================================
\section{Causal Estimands}
\label{w4:sec:estimands}
% ==============================================================

\subsection{The Average Treatment Effect and Its Relatives}

\begin{defn}{ATE, ATT, and ATC}
\label{w4:defn:ate}
For a binary treatment $T \in \{0,1\}$:
\begin{itemize}
  \item The \emph{average treatment effect} (ATE) is
        $\tau_{\mathrm{ATE}} = \E[Y(1) - Y(0)]$.
  \item The \emph{average treatment effect on the treated} (ATT) is
        $\tau_{\mathrm{ATT}} = \E[Y(1) - Y(0) \mid T{=}1]$.
  \item The \emph{average treatment effect on the controls} (ATC) is
        $\tau_{\mathrm{ATC}} = \E[Y(1) - Y(0) \mid T{=}0]$.
\end{itemize}
\end{defn}

The three estimands are related by:
\begin{equation}
  \tau_{\mathrm{ATE}}
  \;=\; P(T{=}1)\,\tau_{\mathrm{ATT}}
       + P(T{=}0)\,\tau_{\mathrm{ATC}}.
  \label{w4:eq:ate-decomp}
\end{equation}
This follows from the law of total expectation applied to $t \in \{0,1\}$.
The ATE and ATT coincide only when the treatment effect does not depend on
who selected into treatment.

None of these quantities is directly observable.  The naive estimator
$\hat\tau_{\mathrm{naive}} = \bar{Y}_{T=1} - \bar{Y}_{T=0}$ estimates
$\E[Y \mid T{=}1] - \E[Y \mid T{=}0] \neq \E[Y(1)] - \E[Y(0)]$,
with the gap being the selection bias induced by endogeneity.

\begin{remark}
ATE, ATT, and ATC coincide when $\E[Y(t) \mid T] = \E[Y(t)]$ for
$t \in \{0,1\}$, i.e.\ when treatment assignment is mean-independent of
potential outcomes.  In a randomized experiment, this holds by design.
\end{remark}

\subsection{Causal Estimands as Functionals of the Structural Equation}
\label{w4:subsec:sem-po}

The potential outcome $Y_i(t)$ was introduced as a \emph{hypothetical}:
the value unit $i$ would have exhibited under treatment $t$.  The SEM
framework of Chapter~1 gives this hypothetical a precise generative
meaning, and thereby expresses the potential-outcome estimands as
functionals of the structural equation for $Y$.

\paragraph{From structural equation to potential outcome.}
In the SEM, the outcome is determined by a structural equation
\begin{equation}
  Y \;=\; g(T,\, \mathbf{X},\, U_Y),
  \label{w4:eq:sem-y}
\end{equation}
where $U_Y$ collects all sources of variation in $Y$ not already
accounted for by $(T, \mathbf{X})$.  The potential outcome under
$\doop(T{=}t)$ is obtained by substituting $t$ for $T$ while holding
everything else fixed:
\begin{equation}
  Y_i(t) \;=\; g(t,\, \mathbf{X}_i,\, U_{Y,i}).
  \label{w4:eq:po-from-sem}
\end{equation}
This is precisely what graph surgery does: the mutilated graph
$\Gcal_{\overline{T}}$ replaces the equation for $T$ with the constant
$t$, leaving the equation for $Y$ unchanged.  Equation~\eqref{w4:eq:po-from-sem}
makes explicit that the potential outcome is a unit-level quantity
determined by $g$, $\mathbf{X}_i$, and the unit's own error $U_{Y,i}$ ---
never by the treatments of other units, which is exactly SUTVA.

\paragraph{ATE and ATT as structural parameters.}
Substituting \eqref{w4:eq:po-from-sem} into the definitions gives:
\begin{align}
  \tau_{\mathrm{ATE}}
    &= \E\!\left[g(1, \mathbf{X}, U_Y) - g(0, \mathbf{X}, U_Y)\right],
    \label{w4:eq:ate-sem} \\
  \tau_{\mathrm{ATT}}
    &= \E\!\left[g(1, \mathbf{X}, U_Y) - g(0, \mathbf{X}, U_Y)
               \mid T{=}1\right].
    \label{w4:eq:att-sem}
\end{align}
Both quantities are averages of the unit-level causal effect
$g(1, \mathbf{X}_i, U_{Y,i}) - g(0, \mathbf{X}_i, U_{Y,i})$
over different reference populations.

\paragraph{The linear SEM as a special case.}
In the Gaussian linear SEM of Chapter~1,
\begin{equation}
  Y \;=\; \beta T + \boldsymbol{\gamma}^{\top}\mathbf{X} + \varepsilon,
  \label{w4:eq:linear-sem}
\end{equation}
the structural equation is additive and separable in $T$, so
$Y_i(t) = \beta t + \boldsymbol{\gamma}^{\top}\mathbf{X}_i + \varepsilon_i$
and the unit-level effect is $Y_i(1) - Y_i(0) = \beta$ for every unit.
Consequently,
\[
  \tau_{\mathrm{ATE}} \;=\; \tau_{\mathrm{ATT}} \;=\; \beta.
\]
The structural coefficient $\beta$ \emph{is} the average treatment effect:
no averaging over heterogeneity is needed because there is none.  This
homogeneity is a special property of the linear additive model, not a
general feature.

\paragraph{Heterogeneous effects.}
In the nonparametric SEM~\eqref{w4:eq:sem-y}, the unit-level effect
$g(1, \mathbf{X}_i, U_{Y,i}) - g(0, \mathbf{X}_i, U_{Y,i})$ varies
across units.  The ATE averages this over the full population; the ATT
averages it over the treated subpopulation.  The two differ whenever
treatment selection correlates with the individual effect size ---
i.e., whenever units who tend to benefit more also tend to self-select
into treatment.

\begin{remark}
The SEM representation~\eqref{w4:eq:po-from-sem} clarifies why the
potential outcomes framework and the do-calculus are \emph{equivalent in
content but different in emphasis}.  The do-calculus works with the
interventional distribution $f(y \mid \doop(T{=}t))$ as a population-level
object and asks when it can be recovered from observational data.  The
potential outcomes framework works with unit-level quantities $Y_i(t)$ and
asks what population summaries (ATE, ATT) are scientifically meaningful.
The SEM ties the two together: $Y_i(t) = g(t, \mathbf{X}_i, U_{Y,i})$
is the unit-level object, and integrating over the distribution of
$(\mathbf{X}, U_Y)$ recovers the interventional distribution.
\end{remark}

% ==============================================================
\section{Ignorability and the Back-Door Criterion}
\label{w4:sec:ignorability}
% ==============================================================

\subsection{The Ignorability Assumption}

The central identifying assumption in observational studies is that
treatment assignment is as good as random after conditioning on observed
covariates $X$.

\begin{defn}{Strong Ignorability \citep{rosenbaum1983central}}
\label{w4:defn:ignorability}
The treatment assignment $T$ is \emph{strongly ignorable} given $X$ if:
\begin{enumerate}
  \item \textbf{Unconfoundedness:}
        $\bigl(Y(0),\, Y(1)\bigr) \;\indep\; T \mid X$.
  \item \textbf{Overlap (positivity):}
        $0 < P(T{=}1 \mid X{=}x) < 1$ for all $x$ in the support of $X$.
\end{enumerate}
\end{defn}

Under strong ignorability, the ATE is identified:
\begin{equation}
  \tau_{\mathrm{ATE}}
  \;=\; \E_X\!\left[\E[Y \mid T{=}1, X] - \E[Y \mid T{=}0, X]\right].
  \label{w4:eq:ate-id}
\end{equation}
This is the back-door adjustment formula of Chapter~3, written in
potential-outcome notation.

\subsection{Ignorability as a Graphical Statement}

\begin{defn}{Faithfulness}
\label{w4:defn:faithfulness}
A distribution $P$ is \textbf{faithful} to a DAG $\Gcal$ if every
conditional independence that holds in $P$ is entailed by
$d$-separation in $\Gcal$: that is, $X \indep Y \mid \mathbf{Z}$ in
$P$ only if $X$ and $Y$ are $d$-separated by $\mathbf{Z}$ in $\Gcal$.

Equivalently, faithfulness rules out \emph{accidental} cancellations of
causal paths: no two directed paths between any pair of nodes exactly
cancel each other's effects for all parameter values.  Under the Markov
property, $d$-separation implies conditional independence; faithfulness
adds the converse, making $d$-separation a complete characterization of
the conditional independence structure of $P$.
\end{defn}

\begin{thm}{Ignorability $\Leftrightarrow$ Back-Door Criterion}
\label{w4:thm:ign-bd}
Under the Markov and faithfulness assumptions for the DAG $\Gcal$,
unconfoundedness $(Y(0), Y(1)) \indep T \mid X$ is equivalent to $X$
satisfying the back-door criterion for the effect of $T$ on $Y$ in $\Gcal$:
\begin{enumerate}
  \item No node in $X$ is a descendant of $T$.
  \item $X$ blocks every back-door path from $T$ to $Y$.
\end{enumerate}
\end{thm}

\noindent\textit{Proof sketch.}
The back-door criterion guarantees that in the mutilated graph
$\Gcal_{\overline{T}}$, conditioning on $X$ renders $Y$ independent of all
paths from $T$ except the causal ones.  In the potential outcome language,
this means the assignment mechanism is mean-independent of $Y(t)$ within
strata of $X$.  The Markov property gives the conditional independence.
\hfill$\square$

\paragraph{Intuition.}
The ignorability condition $(Y(0),Y(1)) \indep T \mid X$ states that,
within strata of $X$, treatment assignment carries no information about
the potential outcomes.  Graphically, every back-door path from $T$ to $Y$
must be blocked by $X$: if such a path remained open, information about
unobserved variables would flow from $T$ to $Y(t)$, violating the
independence.  Thus ignorability is exactly the potential-outcome expression
of the back-door criterion --- the potential outcomes framework
\emph{assumes} the criterion but does not \emph{name} it.

\begin{example}{Labor Training Program}
\label{w4:ex:lalonde}
Following \citet{lalonde1986evaluating} and \citet{dehejia1999causal},
let $T$ = receipt of job training, $Y$ = earnings two years later,
$X$ = (age, education, prior earnings, race, marital status).

The back-door path $T \leftarrow X \to Y$ is blocked by conditioning on
$X$, leaving only the causal path $T \to Y$.  Ignorability holds if this
DAG is correctly specified.  If there is an unobserved variable $U$
(motivation, ability) that affects both training participation and earnings,
$X$ fails the back-door criterion and ignorability fails.
\end{example}

\subsection{Overlap and Positivity}

\begin{warn}{Why Overlap Is Non-Negotiable}
Without overlap, some covariate strata contain only treated or only control
units.  In those strata, $\E[Y \mid T{=}0, X{=}x]$ or
$\E[Y \mid T{=}1, X{=}x]$ is unobservable, and
Equation~\eqref{w4:eq:ate-id} cannot be evaluated at those values of $x$.
Identification of the ATE fails not because the causal structure is wrong,
but because the data do not span the support needed.  The ATT may remain
identified even when full overlap fails, because ATT averages only over the
treated population where $P(T{=}1 \mid X) > 0$ by construction.
\end{warn}

% ==============================================================
\section{Single World Intervention Graphs}
\label{w4:sec:swig}
% ==============================================================

The potential outcomes framework and the do-calculus share a common goal
but differ in representation.  Single World Intervention Graphs (SWIGs),
introduced by \citet{richardson2014ace} and given an accessible practical
treatment by \citet{bezuidenhout2025swigs}, provide a unified graphical
representation that encodes both.  A SWIG is constructed directly from a
DAG and is therefore easy to adopt for anyone familiar with the graphical
language of Chapters~2--3.  Its key advantage is making potential outcomes
and identifiability assumptions \emph{visually explicit} in a single
diagram, rather than leaving them implicit as in a standard DAG.

\subsection{The SWIG Construction}
\label{w4:subsec:swig-def}

In the original DAG $\Gcal$, the node $T$ represents the treatment as it
naturally occurs.  Under $\doop(T{=}t)$, however, the treatment is
externally fixed.  The SWIG separates these two roles by \emph{splitting}
$T$ into two pieces displayed side by side in the graph:
\begin{itemize}
  \item a \textbf{random half} (left side, capital letter $T$): represents
        the treatment value that would naturally occur, and retains all
        incoming edges from $T$'s causal parents;
  \item a \textbf{fixed half} (right side, lowercase letter $t$): represents
        the externally imposed intervention value, retains all outgoing
        edges that formerly left $T$, and has \emph{no} direct edge to or
        from the random half.
\end{itemize}

\begin{defn}{Single World Intervention Graph \citep{richardson2014ace}}
\label{w4:defn:swig}
The \emph{SWIG} $\Gcal(t)$ is constructed from $\Gcal$ by:
\begin{enumerate}
  \item Replacing $T$ with two nodes: random half $T$ (all incoming edges
        retained) and fixed half $t$ (all outgoing edges retained).
  \item Severing the direct connection between the two halves.
  \item Redrawing causal arrows at the split node:
        \emph{all arrows into the split node land on the random half};
        \emph{all arrows departing the split node leave from the fixed half}.
  \item Relabeling every descendant $V$ of $T$ in $\Gcal$ as $V(t)$,
        indicating it is a potential outcome under $\doop(T{=}t)$.
\end{enumerate}
\end{defn}

Figure~\ref{w4:fig:swig} illustrates the construction for the simple
confounded-treatment graph.  Splitting $A$ into its random and fixed
halves immediately reveals $Y(a)$ as a potential outcome and makes the
key independence $A \indep Y(a) \mid L$ graphically transparent.

\begin{figure}[ht]
\centering
\begin{tikzpicture}[>=stealth]
  \tikzstyle{rnd}=[circle,draw=isublue,fill=defbg,thick,
                   minimum size=9mm,font=\small\bfseries]
  \tikzstyle{fix}=[rectangle,draw=isublue,fill=gray!15,thick,
                   minimum size=9mm,font=\small\bfseries]
  %--- DAG (left panel) ---
  \begin{scope}[xshift=0cm]
    \node[rnd] (L)  at (1.4, 1.5) {$L$};
    \node[rnd] (A)  at (0,   0)   {$A$};
    \node[rnd] (Y)  at (2.8, 0)   {$Y$};
    \draw[->,thick,color=isublue](L)--(A);
    \draw[->,thick,color=isublue](L)--(Y);
    \draw[->,thick,color=isublue](A)--(Y);
    \node[font=\small\itshape\color{darkgrey}] at (1.4,-0.85) {DAG};
  \end{scope}
  %--- Arrow ---
  \node[font=\large] at (3.9, 0.4) {$\Rightarrow$};
  %--- SWIG (right panel) ---
  \begin{scope}[xshift=4.9cm]
    \node[rnd] (Lr)  at (1.9, 1.5) {$L$};
    \node[rnd] (Ar)  at (0.4, 0)   {$A$};
    \node[fix] (af)  at (1.5, 0)   {$a$};
    \node[rnd] (Ya)  at (3.5, 0)   {$Y(a)$};
    \draw[->,thick,color=isublue](Lr)--(Ar);
    \draw[->,thick,color=isublue](Lr)--(Ya);
    \draw[->,thick,color=isublue](af)--(Ya);
    \node[font=\small\itshape\color{darkgrey}] at (1.9,-0.85)
      {SWIG $\Gcal(a)$};
  \end{scope}
\end{tikzpicture}
\caption{Converting the confounded-treatment DAG ($L$ is a confounder)
  into the SWIG for $\doop(A{=}a)$.  The treatment node splits: the
  random half $A$ (circle) retains the incoming arrow from $L$; the
  fixed half $a$ (square) carries the outgoing arrow to $Y(a)$.
  The potential outcome $Y(a)$ is caused by $a$, not by the natural
  treatment $A$.  D-separation in the SWIG gives
  $A \indep Y(a) \mid L$ (conditional ignorability).}
\label{w4:fig:swig}
\end{figure}

\subsection{SWIG D-Separation Rules}
\label{w4:subsec:swig-dsep}

All standard d-separation rules from Chapter~2 apply within a SWIG.
In addition, the split-node structure introduces two rules specific to
fixed (intervention) nodes \citep{richardson2014ace,bezuidenhout2025swigs}.

\begin{tcolorbox}[defstyle, title={SWIG-Specific D-Separation Rules}]
\begin{enumerate}
  \item \textbf{Paths through a fixed node are always blocked.}
        A path that \emph{passes through} the fixed half (e.g.\
        $A \to a \to Y(a)$) is always blocked, because in the single
        world where everyone receives treatment $a$, the natural value
        $A$ cannot affect $Y(a)$ through the fixed half.  Note: the
        path \emph{originating} from $a$ (i.e.\ $a \to Y(a)$) is
        \emph{not} blocked --- it is a direct causal arrow and remains
        open.
  \item \textbf{A fixed node connects to a random node only via open
        paths.}  A fixed node can be d-connected to a random node only
        if the path between them is not already blocked by the standard
        d-separation rules.
\end{enumerate}
\end{tcolorbox}

\paragraph{Key consequence: recovering ignorability from the SWIG.}
In Figure~\ref{w4:fig:swig}, the two paths from the random half $A$ to
$Y(a)$ are:
\begin{itemize}
  \item $A \to a \to Y(a)$: blocked by Rule~1 (transits the fixed half).
  \item $A \leftarrow L \to Y(a)$: open unless $L$ is conditioned on.
\end{itemize}
Conditioning on $L$ blocks the second path.  D-separation in the SWIG
then gives
\[
  A \;\indep\; Y(a) \;\big|\; L,
\]
which is exactly the \emph{conditional ignorability} assumption of
Definition~\ref{w4:defn:ignorability}, now \emph{derived as a
graph-theoretic consequence} from the SWIG rather than asserted as a
modeling assumption.

\begin{remark}
More than one node may be split in a single SWIG.  In longitudinal
settings with time-varying treatments $(A_0, A_1)$, splitting both nodes
yields the sequential-randomization SWIG, from which the two
\emph{sequential exchangeability} conditions needed to identify the
g-formula follow directly by d-separation.  This construction is
developed in detail when longitudinal treatments are covered in a later chapter.
\end{remark}

\subsection{What SWIGs Achieve}
\label{w4:subsec:swig-achieve}

\paragraph{1.\ Potential outcomes as random variables.}
In $\Gcal(t)$, the variable $Y(t)$ has well-defined parents $t$ and $U$.
Its distribution is the interventional distribution
$f(y \mid \doop(T{=}t))$, establishing
Proposition~\ref{w4:prop:po-do} graphically.

\paragraph{2.\ The mutilated graph as a special case.}
Marginalizing out the random half $T$ from $\Gcal(t)$ yields exactly the
mutilated graph $\Gcal_{\overline{T}}$: the SWIG is a more detailed version
that tracks what would have determined $T$ in the absence of intervention.

\paragraph{3.\ Cross-world independence made visible.}
Unconfoundedness $(Y(t) \indep T \mid X)$ is a d-separation statement
\emph{in the SWIG} $\Gcal(t)$, verifiable by inspecting the graph.

\begin{warn}{Cross-World Independence and Its Limits}
Consider $Y(1) \indep Y(0)$: $Y(1)$ lives in SWIG $\Gcal(1)$ and $Y(0)$
lives in SWIG $\Gcal(0)$.  No unit is ever observed in both worlds
simultaneously, so no data can directly test this independence.  More
generally, statements such as $Y(t) \indep T(t') \mid X$ for $t \neq t'$
involve variables from \emph{different SWIGs}.  SWIGs make such
assumptions explicit: the analyst must specify a joint distribution across
multiple SWIGs --- a commitment that potential-outcome notation alone does
not force.
\end{warn}

\subsection{Ignorability in the SWIG}
\label{w4:subsec:swig-ign}

In $\Gcal(t)$, the random half $T$ retains all its natural parents.
The condition $(Y(t) \indep T \mid X)_{\Gcal(t)}$ holds if and only if
$X$ blocks all paths between $T$ and $Y(t)$ in $\Gcal(t)$.  Since $Y(t)$
descends from the fixed half $t$ (not from the random half $T$), such paths
must pass through $X$ --- which is exactly the back-door condition of
Theorem~\ref{w4:thm:ign-bd}, now expressed as a single d-separation
statement.  The SWIG thus provides a \emph{unified} perspective:
ignorability is neither an assumption of the potential outcomes framework
nor an assumption of the do-calculus in isolation, but a graph-structural
property of the SWIG that both frameworks share.

% ==============================================================
\section{Where the Frameworks Agree and Diverge}
\label{w4:sec:comparison}
% ==============================================================

\subsection{A Systematic Comparison}

\begin{center}
\renewcommand{\arraystretch}{1.5}
\begin{tabular}{@{}p{3.2cm}p{3.4cm}p{3.4cm}p{3.4cm}@{}}
\toprule
\textbf{Task}
  & \textbf{Potential Outcomes}
  & \textbf{Do-Calculus / DAG}
  & \textbf{SEM} \\
\midrule
Define causal estimands (ATE, ATT)
  & \checkmark\ Native language
  & Via $\E[Y \mid \doop(t)]$
  & Via structural equations \\
\addlinespace
Encode causal assumptions
  & Stated informally; graph hidden
  & \checkmark\ Explicit directed edges
  & \checkmark\ Structural equations \\
\addlinespace
Read conditional independence
  & Not possible without a graph
  & \checkmark\ d-separation
  & With graph only \\
\addlinespace
Identification from observational data
  & Via ignorability assumptions
  & \checkmark\ Back-door, front-door, do-calculus
  & \checkmark\ Via exclusion restrictions \\
\addlinespace
Likelihood construction
  & Moment restrictions only
  & \checkmark\ Via back-door/front-door
  & \checkmark\ Via structural model \\
\addlinespace
Cross-world assumptions (e.g.\ monotonicity)
  & \checkmark\ Natural to state
  & Requires multiple SWIGs
  & Implicit in structural model \\
\addlinespace
Standard in statistics/epidemiology
  & \checkmark\ Dominant
  & Growing rapidly
  & Econometrics \\
\bottomrule
\end{tabular}
\end{center}

\begin{tcolorbox}[defstyle, title={Our Position in This Course}]
Both frameworks are indispensable.  We use \emph{potential outcomes} to
\emph{define} causal estimands.  We use the \emph{do-calculus and DAGs} for
\emph{identification}.  SWIGs provide the formal bridge when both are needed
simultaneously.

The do-operator language is preferred throughout this course because it
makes the interventional/observational distinction \emph{syntactically
enforced}: $f(y \mid \doop(t))$ cannot be confused with $f(y \mid T{=}t)$,
whereas $\E[Y(t)]$ carries no such syntactic marker.  For likelihood
construction in Part~III, this distinction determines whether the
estimating equations are correctly derived.
\end{tcolorbox}

% ==============================================================
\section{Part I Synthesis: The Gaussian Linear IV Model}
\label{w4:sec:synthesis}
% ==============================================================

The Gaussian linear IV model, introduced in Chapter~1, can now be given a
complete account in all three languages.  Table~\ref{w4:tab:synthesis}
summarizes the translation.

\begin{table}[h]
\centering
\small
\renewcommand{\arraystretch}{1.5}
\begin{tabular}{@{}p{2.6cm}p{3.6cm}p{4.2cm}p{3.6cm}@{}}
\toprule
\textbf{Concept}
  & \textbf{SEM}
  & \textbf{DAG / Do-Calculus}
  & \textbf{Potential Outcomes} \\
\midrule
Treatment model
  & $T = \pi Z + \delta$,\; $\delta \sim \mathcal{N}(0,\tau^2)$
  & $Z\to T\to Y$; hidden $U\to T$, $U\to Y$; graph surgery yields $\Gcal_{\overline{T}}$
  & $T(z) = \pi z + \delta$; potential treatment indexed by $z$ \\
Outcome model
  & $Y = \beta T + \varepsilon$,\; $\mathrm{Cov}(\varepsilon,\delta)\neq 0$
  & $f(y\mid\doop(T{=}t))$ is the causal quantity; OLS targets $f(y\mid T{=}t)$
  & $Y(t) = \beta t + \varepsilon$; linear homogeneous potential outcome \\
Relevance
  & $\pi \neq 0$; $Z$ predicts $T$
  & $Z \nindep T$ in $\Gcal$; $Z$ and $T$ are d-connected
  & $\E[T(1)] \neq \E[T(0)]$ \\
Exclusion restriction
  & $Z$ absent from outcome equation
  & No direct edge $Z\to Y$;\; $f(y\mid\doop(T{=}t),z) = f(y\mid\doop(T{=}t))$
  & $Y_i(t, z) = Y_i(t, z')$ for all $z, z'$ \\
Exogeneity
  & $\mathrm{Cov}(Z,\delta) = \mathrm{Cov}(Z,\varepsilon) = 0$
  & $Z \indep U$ in $\Gcal$; no back-door path $Z\leftarrow\cdots\to Y$
  & $Z \indep (Y(0), Y(1), T(0), T(1))$ \\
Identified estimand
  & $\beta = \mathrm{Cov}(Y,Z)/\mathrm{Cov}(T,Z)$ (Wald)
  & $\E[Y\mid\doop(T{=}1)] - \E[Y\mid\doop(T{=}0)]$; IV licenses graph surgery
  & $\beta = \tau_{\mathrm{LATE}} = \tau_{\mathrm{ATE}}$ (homogeneous effects) \\
Endogeneity bias
  & $\mathrm{plim}\;\hat\beta_{\mathrm{OLS}} = \beta + \rho\sigma/\tau$
  & Back-door path $T\leftarrow U\to Y$ open; do-calculus blocks it via $Z$
  & Selection bias: $\E[Y\mid T{=}1] - \E[Y\mid T{=}0] \neq \tau_{\mathrm{ATE}}$ \\
\bottomrule
\end{tabular}
\caption{The Gaussian linear IV model in all three languages.  In the
linear homogeneous model $\tau_{\mathrm{LATE}} = \tau_{\mathrm{ATE}} = \beta$,
so all three languages identify the same quantity.}
\label{w4:tab:synthesis}
\end{table}

\paragraph{Why the Homogeneous Linear Case is Special.}
In this model, $Y(t) = \beta t + \varepsilon$ with the \emph{same}
coefficient $\beta$ for all units.  Consequently:
\[
  \tau_{\mathrm{LATE}}
  = \E[Y(1) - Y(0) \mid \text{compliers}]
  = \beta
  = \E[Y(1) - Y(0)]
  = \tau_{\mathrm{ATE}}.
\]
This collapse is a special feature of the linear homogeneous model.  In the
binary treatment, heterogeneous-effects setting of Chapter~7, the distinction
between LATE and ATE is fundamental.

\paragraph{Part I in Retrospect.}

\begin{center}
\begin{tikzpicture}[
    box/.style = {rectangle, rounded corners=5pt, draw=isublue, fill=defbg,
                  thick, text width=6.8cm, align=left, minimum height=1.0cm,
                  inner sep=8pt, font=\small},
    arr/.style = {-{Stealth[length=7pt,width=5pt]}, line width=1.6pt,
                  color=accent},
    lbl/.style = {font=\scriptsize\bfseries\color{isublue},
                  minimum width=1.4cm, align=center},
    node distance = 0.4cm
  ]
  \node[box] (C1) {\textbf{Chapter 1.} Interventional vs.\ observational:
    $f(y \mid \doop(t)) \neq f(y \mid T{=}t)$.  The gap is endogeneity bias.};
  \node[box, below=0.55cm of C1] (C2)
    {\textbf{Chapter 2.} DAGs and d-separation: read independence from the
    graph; identify which paths are confounded.};
  \node[box, below=0.55cm of C2] (C3)
    {\textbf{Chapter 3.} Do-calculus: translate graph surgery into
    identification formulas (back-door, front-door).};
  \node[box, below=0.55cm of C3] (C4)
    {\textbf{Chapter 4.} Potential outcomes: define the estimands
    ($\tau_{\mathrm{ATE}}, \tau_{\mathrm{ATT}}$); connect them to the
    do-calculus and to the structural equation.  SWIGs unify all three
    languages in a single graph.};
  \draw[arr] (C1) -- (C2);
  \draw[arr] (C2) -- (C3);
  \draw[arr] (C3) -- (C4);
  \node[lbl, left=0.6cm of C1] {Ch.\ 1};
  \node[lbl, left=0.6cm of C2] {Ch.\ 2};
  \node[lbl, left=0.6cm of C3] {Ch.\ 3};
  \node[lbl, left=0.6cm of C4] {Ch.\ 4};
\end{tikzpicture}
\end{center}

\medskip
Part~II (Chapters~5--8) applies this foundation to randomization and
back-door adjustment (Chapter~5), propensity score methods (Chapter~6),
instrumental variables (Chapter~7), and mediation analysis (Chapter~8).

% ==============================================================
\section{Summary}
\label{w4:sec:summary}
% ==============================================================

\begin{enumerate}
  \item The potential outcome $Y(t)$ is the outcome that would be observed
        under the intervention $T = t$.  Under SUTVA and consistency,
        $Y_i = Y_i(T_i)$ and $Y(t) \sim f(y \mid \doop(T{=}t))$.

  \item The ATE and ATT are averages of the unit-level causal effect
        $Y_i(1) - Y_i(0) = g(1, \mathbf{X}_i, U_{Y,i}) - g(0, \mathbf{X}_i, U_{Y,i})$
        over the full population and the treated subpopulation respectively.
        In the linear SEM, both equal the structural coefficient $\beta$.
        They diverge under heterogeneous effects when treatment selection
        correlates with individual effect size.

  \item Strong ignorability $(Y(0),Y(1)) \indep T \mid X$ with overlap
        identifies the ATE via Equation~\eqref{w4:eq:ate-id}.  By
        Theorem~\ref{w4:thm:ign-bd}, unconfoundedness is equivalent to the
        back-door criterion.

  \item SWIGs \citep{richardson2014ace,bezuidenhout2025swigs} are constructed
        from a DAG by splitting each intervened node into a random half
        (capital letter, incoming edges) and a fixed half (lowercase letter,
        outgoing edges), and relabeling descendants as potential outcomes.
        Two SWIG-specific d-separation rules supplement the standard ones:
        (i)~paths \emph{transiting} a fixed node are always blocked;
        (ii)~a fixed node connects to a random node only via paths that are
        open under the standard rules.  Together these rules make ignorability
        readable directly from the graph.  In longitudinal settings, splitting
        multiple treatment nodes yields sequential exchangeability conditions
        by the same d-separation logic (to be developed in a later chapter).

  \item The frameworks are complementary: potential outcomes define
        estimands; the do-calculus identifies them.  SWIGs provide the
        formal bridge when both are needed simultaneously.
\end{enumerate}

\subsection*{From Ignorability to Randomization}

In observational studies, ignorability must be justified by substantive
knowledge encoded in a causal graph.  Because unobserved confounding can
never be ruled out empirically, this justification is often debated.

Randomized experiments provide a fundamentally different solution.  When
treatment is assigned randomly, $(Y(0),Y(1)) \indep T$ holds by design,
guaranteeing ignorability \emph{without any covariate adjustment}.  Chapter~5
studies randomized experiments as the canonical design where causal effects
are identifiable directly from the data-generating mechanism.

\bigskip
\begin{center}
  \large\textbf{Causal inference}
  $=$ counterfactual questions
  $+$ graphical assumptions
  $+$ statistical estimation.
\end{center}

\subsection*{Identification vs.\ Estimation}

The first four chapters have focused on \emph{identification}: whether a
causal quantity $\tau = \E[Y(1) - Y(0)]$ can be written as
$\tau = \Phi(P(Y,T,X))$ for some functional $\Phi$ of the observed
distribution.  Estimation --- constructing $\hat\tau$ from a finite sample
$(Y_i, T_i, X_i)_{i=1}^n$ to approximate $\Phi(P)$ --- is studied in
Part~III, covering regression adjustment, inverse probability weighting,
doubly robust estimators, and instrumental variables.

% ==============================================================
\newpage
\addcontentsline{toc}{section}{Problems}
\section*{Problems}
% ==============================================================

\begin{enumerate}[label=\textbf{\arabic*.}, leftmargin=*, itemsep=10pt]

  \item \textbf{SUTVA and consistency.}
  Suppose $n = 3$ units receive binary treatments $(T_1, T_2, T_3)$ and
  unit $i$'s outcome may depend on all three treatments:
  $Y_i = Y_i(T_1, T_2, T_3)$.
  \begin{enumerate}[label=(\alph*)]
    \item How many potential outcomes does unit~1 have?  Write them out
          explicitly.
    \item SUTVA imposes no interference: $Y_i(T_1,T_2,T_3) = Y_i(T_i)$.
          How many distinct potential outcomes remain?
    \item Give a real-world example where no-interference plausibly holds
          and one where it plausibly fails.  In the latter case, explain
          what identification strategy (if any) remains available.
  \end{enumerate}

  \item \textbf{ATE, ATT, and selection bias.}
  Let $(Y(0), Y(1), T) \sim P$ with $P(T{=}1) = 0.5$.  Suppose
  $\E[Y(1)] = 3$, $\E[Y(0)] = 1$, $\E[Y(1) \mid T{=}1] = 4$,
  $\E[Y(0) \mid T{=}1] = 2$, $\E[Y(1) \mid T{=}0] = 2$,
  $\E[Y(0) \mid T{=}0] = 0$.
  \begin{enumerate}[label=(\alph*)]
    \item Compute the ATE and ATT.  Are they equal?
    \item Compute the naive estimator $\E[Y \mid T{=}1] - \E[Y \mid T{=}0]$.
          Decompose the gap between this and the ATE into a selection bias
          term and an ATT--ATE difference.
    \item Under what graphical condition on the DAG would the naive
          estimator equal the ATE?  State this condition in both potential
          outcome language (ignorability) and do-calculus language
          (back-door criterion).
  \end{enumerate}

  \item \textbf{Causal estimands from the structural equation.}
  Consider the nonparametric SEM $Y = g(T, X, U_Y)$ with binary treatment
  $T \in \{0,1\}$, observed covariate $X$, and unobserved error $U_Y$
  independent of $(T, X)$ in the mutilated graph.
  \begin{enumerate}[label=(\alph*)]
    \item Write $\tau_{\mathrm{ATE}}$ and $\tau_{\mathrm{ATT}}$ as
          expectations of $g(1, X, U_Y) - g(0, X, U_Y)$ over the
          appropriate distribution.  Under what condition do they coincide?
    \item Specialize to the linear SEM $g(t, x, u) = \alpha + \beta t
          + \gamma x + u$.  Show that the unit-level causal effect
          $Y_i(1) - Y_i(0)$ is constant across all units, and hence
          $\tau_{\mathrm{ATE}} = \tau_{\mathrm{ATT}} = \beta$.
    \item Now consider the heterogeneous SEM $g(t, x, u) = (\alpha + u)\,t
          + \gamma x$, where $U_Y \sim \mathcal{N}(0, \sigma^2)$ and
          $\mathrm{Cov}(U_Y, T) = \rho$.  Compute $\tau_{\mathrm{ATE}}$
          and $\tau_{\mathrm{ATT}}$.  Show that they differ when $\rho \neq 0$,
          and interpret this difference in words.
    \item In part~(c), explain why OLS estimation of $Y$ on $T$ and $X$
          does not recover either $\tau_{\mathrm{ATE}}$ or
          $\tau_{\mathrm{ATT}}$ when $\rho \neq 0$.  What additional
          structure would be needed for identification?
  \end{enumerate}

  \item \textbf{SWIGs and ignorability.}
  Consider the DAG: $X \to T$, $X \to Y$, $T \to Y$, with $X$ fully
  observed.
  \begin{enumerate}[label=(\alph*)]
    \item Construct the SWIG $\Gcal(t)$ by splitting $T$ into its random
          and fixed halves (Definition~\ref{w4:defn:swig}).  Draw the result,
          labeling the random half, fixed half, and the potential outcome $Y(t)$.
    \item In $\Gcal(t)$, identify all paths between the random half $T$
          and $Y(t)$.  Apply the SWIG d-separation rules to determine
          which paths are blocked and which are open before conditioning.
    \item Use d-separation in $\Gcal(t)$ to verify that
          $(Y(t) \indep T \mid X)_{\Gcal(t)}$ holds.  Conclude that $X$
          satisfies the back-door criterion, confirming
          Theorem~\ref{w4:thm:ign-bd}.
    \item Now add a hidden common cause $U \to T$, $U \to Y$.  Draw the
          revised SWIG.  Does the ignorability argument still hold?
          State the correct conclusion and identify what additional
          structure (if any) would be needed for identification.
  \end{enumerate}

\end{enumerate}

% ==============================================================
% Bibliography note: the following entry must be added to
% causal_inference.bib if not already present.
%
% @article{bezuidenhout2025swigs,
%   author  = {Bezuidenhout, Dana and Forthal, Sarah and
%              Rudolph, Kara and Lamb, Matthew R.},
%   title   = {Single world intervention graphs ({SWIGs}):
%              a practical guide},
%   journal = {American Journal of Epidemiology},
%   year    = {2025},
%   volume  = {194},
%   pages   = {2047--2052},
%   doi     = {10.1093/aje/kwae353}
% }
% ==============================================================

\ifSubfilesClassLoaded{
  \newpage
  \section*{References}
  \addcontentsline{toc}{section}{References}
  \bibliographystyle{plainnat}
  \bibliography{causal_inference}
}{}

\end{document}

